\chapter{Analyse des besoins et étude de l'existant}
\label{chap:analyse-besoins-existant}
\begin{spacing}{1.1}
\minitoc
\thispagestyle{MyStyle}
\end{spacing}
\newpage

\section{Introduction}
Dans la continuité du cadre général présenté au chapitre \ref{chap:cadre-general}, ce chapitre précise les besoins à satisfaire et positionne la solution proposée par rapport à l'existant. L'objectif est double: formaliser des exigences exploitables pour la conception, puis justifier les choix techniques au regard des limites observées dans les approches actuelles.

\section{Analyse des besoins}
\subsection{Besoins fonctionnels}
Les besoins fonctionnels retenus couvrent l'ensemble du cycle d'enrichissement des données produit:
\begin{itemize}
  \item Déclencher l'enrichissement automatique d'une fiche produit à partir d'images, d'attributs existants et de descriptions disponibles;
  \item Générer un ensemble d'attributs structurés (couleur dominante et secondaire, matière, type/catégorie, motif, saisonnalité, description structurée);
  \item Affecter un score de confiance par attribut afin d'aider la décision métier;
  \item Permettre la validation, la correction ou le rejet des attributs proposés via une interface de pilotage;
  \item Tracer les décisions de validation (audit trail: utilisateur, date, valeur initiale, valeur finale);
  \item Injecter les attributs validés dans SAP Hybris avec gestion des statuts et mécanismes de reprise en cas d'erreur;
  \item Réindexer les données enrichies dans Apache Solr pour améliorer filtres, facettes et pertinence de recherche;
  \item Exposer les données enrichies via API REST pour la réutilisation par d'autres composants.
\end{itemize}

\subsection{Besoins non fonctionnels}
Au-delà des fonctionnalités, plusieurs exigences de qualité structurent la solution:
\begin{itemize}
  \item \textbf{Performance}: traitement en mode batch compatible avec le volume catalogue (jusqu'à plusieurs milliers de produits par lot);
  \item \textbf{Fiabilité}: gestion robuste des erreurs d'API (retry, journalisation, reprise);
  \item \textbf{Auditabilité}: historisation des modifications pour garantir la traçabilité métier;
  \item \textbf{Maintenabilité}: ajout de nouveaux attributs sans refonte majeure du pipeline;
  \item \textbf{Interopérabilité}: alignement avec l'écosystème Java/Spring, Hybris et Solr;
  \item \textbf{Évolutivité}: capacité à intégrer ultérieurement des fonctionnalités multilingues et des contrôles de non-conformité.
\end{itemize}

\subsection{Acteurs et cas d'usage}
La Figure \ref{fig:usecase-global} synthétise les interactions majeures entre les acteurs métier et la plateforme. Même si le schéma est générique, il met en évidence une idée centrale: chaque action sensible est conditionnée par l'authentification et s'inscrit dans une logique de gouvernance des données.

\begin{figure}[H]
  \centering
  \includegraphics[width=0.95\linewidth]{diag1.jpg}
  \caption{Diagramme de cas d'usage global de la plateforme}
  \label{fig:usecase-global}
\end{figure}

En transposant ce principe au périmètre du PFE, les cas d'usage prioritaires sont:
\begin{itemize}
  \item UC1: Déclencher l'enrichissement d'un produit;
  \item UC2: Valider / corriger / rejeter les attributs générés;
  \item UC3: Importer les attributs validés dans Hybris;
  \item UC4: Réindexer les produits concernés dans Solr;
  \item UC5: Suivre les indicateurs de qualité des enrichissements.
\end{itemize}

\subsection{Priorisation des user stories}
La planification fonctionnelle initiale est résumée dans la Figure \ref{fig:planning-user-stories}. Elle met en évidence un découpage progressif: d'abord les fonctions essentielles, puis les fonctions d'administration et de pilotage.

\begin{figure}[H]
  \centering
  \includegraphics[width=0.95\linewidth]{tabm.drawio.png}
  \caption{Priorisation initiale des user stories par release}
  \label{fig:planning-user-stories}
\end{figure}

\section{Étude de l'existant}
\subsection{Objectif et protocole des tests}
L'étude de l'existant ne se limite pas à une comparaison théorique. Elle s'appuie sur des essais pratiques réalisés sur des cas réels de fiches produit (texte, métadonnées, images multiples), afin de répondre à trois questions:
\begin{itemize}
  \item les solutions \textbf{gratuites locales} sont-elles exploitables dans nos contraintes matérielles;
  \item les \textbf{APIs gratuites} sont-elles suffisantes pour un flux e-commerce réel;
  \item les \textbf{APIs payantes} apportent-elles le niveau de robustesse attendu pour industrialiser la solution.
\end{itemize}

Les tests ont été menés sur des scénarios de difficulté croissante: texte seul, image unique, multi-images, puis images + métadonnées.

\subsection{Axe 1 -- Solutions gratuites locales (Ollama)}
\subsubsection{Résultats observés}
Les essais locaux avec des modèles multimodaux open-source (Gemma, Qwen-VL, Moondream) montrent une limite structurelle de l'infrastructure disponible:
\begin{itemize}
  \item latence très élevée (plusieurs minutes même pour une réponse simple);
  \item forte pression CPU/RAM en l'absence de GPU dédié;
  \item instabilités ou échecs sur les cas multi-images.
\end{itemize}

\begin{figure}[H]
  \centering
  \includegraphics[width=0.9\linewidth]{image.png}
  \caption{Exemple d'échec et de surcharge en exécution locale}
  \label{fig:local-ollama-error}
\end{figure}

\begin{figure}[H]
  \centering
  \includegraphics[width=0.9\linewidth]{image_1.png}
  \caption{Saturation des ressources machine en test local}
  \label{fig:local-ollama-cpu}
\end{figure}

\subsubsection{Variante Colab GPU gratuit}
Pour contourner la limite matérielle locale, des tests ont été réalisés sur Google Colab (GPU gratuit). Les modèles lourds restent néanmoins difficiles à exploiter opérationnellement: temps de réponse long, consommation GPU proche de la limite mémoire, et faible marge de stabilité pour des traitements en lot.

\begin{figure}[H]
  \centering
  \includegraphics[width=0.95\linewidth]{test_qwen3.5.27b_avec_colab.png}
  \caption{Test Qwen 3.5 27B sur Colab GPU gratuit}
  \label{fig:colab-qwen}
\end{figure}

\begin{figure}[H]
  \centering
  \includegraphics[width=0.95\linewidth]{test_gemma31b_avec_colab.png}
  \caption{Test Gemma 31B sur Colab GPU gratuit}
  \label{fig:colab-gemma}
\end{figure}

\textbf{Verdict de l'axe 1:} viable pour exploration et POC technique, non viable pour un pipeline de production stable dans notre contexte.

\subsection{Axe 2 -- APIs gratuites (OpenRouter, Groq, AI Studio Free)}
\subsubsection{OpenRouter (free tier)}
Les tests OpenRouter confirment un comportement acceptable sur des entrées simples (texte seul, parfois image unique), mais les cas cibles du projet (multi-images + métadonnées) provoquent régulièrement:
\begin{itemize}
  \item \textbf{context overflow};
  \item \textbf{HTTP 429} (Too Many Requests);
  \item dégradation de robustesse entre appels successifs.
\end{itemize}

\begin{figure}[H]
  \centering
  \includegraphics[width=0.95\linewidth]{image_2.png}
  \caption{Exemple de limite contextuelle sur OpenRouter en multi-images}
  \label{fig:openrouter-overflow}
\end{figure}

\subsubsection{Groq (free tier)}
Sur Groq, la latence est bonne, mais la contrainte TPM (tokens par minute) devient bloquante sur un produit enrichi avec images et métadonnées. Les erreurs \texttt{HTTP 413} et \texttt{rate\_limit\_exceeded} empêchent la continuité du flux.

\begin{figure}[H]
  \centering
  \includegraphics[width=0.95\linewidth]{image_3.png}
  \caption{Erreur de quota tokens observée sur Groq}
  \label{fig:groq-413}
\end{figure}

\subsubsection{Google AI Studio (gratuit)}
Les essais montrent une meilleure qualité de réponse que plusieurs alternatives gratuites, mais les mêmes limites restent présentes: quotas gratuits serrés, erreurs 429 sous charge, et absence de garantie de service.

\begin{figure}[H]
  \centering
  \includegraphics[width=0.95\linewidth]{image_4.png}
  \caption{Exemple d'utilisation AI Studio avec contrainte de ressources gratuites}
  \label{fig:aistudio-free}
\end{figure}

\textbf{Verdict de l'axe 2:} utile pour prototypage rapide, insuffisant pour une chaîne fiable d'enrichissement e-commerce.

\subsection{Axe 3 -- APIs payantes (candidats de production)}
Les résultats des axes 1 et 2 conduisent à considérer des APIs payantes pour la phase d'industrialisation. Le Tableau \ref{tab:comparaison-payant} synthétise les options les plus crédibles.

\begin{table}[H]
  \centering
  \caption{Comparaison des APIs payantes candidates pour la production}
  \label{tab:comparaison-payant}
  \begin{tabular}{|p{2.6cm}|p{3.1cm}|p{2.1cm}|p{2.0cm}|p{3.2cm}|}
    \hline
    \textbf{Solution} & \textbf{Forces} & \textbf{Latence} & \textbf{Coût} & \textbf{Risque principal} \\
    \hline
    GPT-4o (OpenAI) & Très bonne qualité multimodale, sortie JSON robuste & Faible à moyenne & Élevé & Coût à l'échelle et gouvernance des données \\
    \hline
    Gemini 2.5 (Google) & Bonne qualité visuelle, contexte large & Moyenne & Moyen à élevé & Stabilité des quotas selon plan choisi \\
    \hline
    Claude Sonnet (Anthropic) & Excellente compréhension sémantique, bonne rigueur de sortie & Faible à moyenne & Élevé & Coût et intégration selon stack retenue \\
    \hline
  \end{tabular}
\end{table}

Sur le plan métier, ces offres payantes répondent mieux aux exigences critiques du projet: robustesse, continuité de service, meilleur support des scénarios multi-images et gestion contrôlée des quotas.

\subsection{Synthèse de décision}
L'étude expérimentale met en évidence une trajectoire claire:
\begin{itemize}
  \item \textbf{Gratuit local}: adapté à la recherche et aux démonstrations ponctuelles, mais trop instable pour un flux opérationnel;
  \item \textbf{APIs gratuites}: utiles pour explorer rapidement des modèles, mais contraintes de quotas incompatibles avec un usage soutenu;
  \item \textbf{APIs payantes}: option la plus réaliste pour atteindre la robustesse attendue en environnement e-commerce.
\end{itemize}

La stratégie recommandée est donc hybride: utiliser les offres gratuites pour le prototypage initial, puis basculer vers une API payante pour la phase de validation industrielle et de déploiement progressif.

\section{Conclusion}
Ce chapitre a transformé le besoin initial en exigences structurées et mesurables, puis a démontré les limites de l'existant et l'intérêt d'une approche hybride orientée IA multimodale. Les éléments établis ici serviront de base à la spécification détaillée et à la conception d'architecture du chapitre suivant.
